% ─────────────────────────────────────────────────────────────────────
%  A Layered, Axiom-Isolated Lean 4 Formalization of the Congruent
%  Number Problem up to Tunnell's Criterion
%
%  Manuscript prepared for submission to
%  Annals of Formalized Mathematics (AFM) — April 2026 (R1)
%  Repository: https://github.com/papanokechi/congruent-relay
% ─────────────────────────────────────────────────────────────────────
\documentclass[11pt,a4paper]{article}

% ── Packages ──────────────────────────────────────────────────────────
\usepackage[utf8]{inputenc}
\usepackage[T1]{fontenc}
\usepackage{lmodern}
\usepackage{amsmath,amssymb,amsthm}
\usepackage{mathtools}
\usepackage{booktabs}
\usepackage{enumitem}
\usepackage{hyperref}
\usepackage[margin=1in]{geometry}
\usepackage{xcolor}
\usepackage{listings}
\usepackage{microtype}
\usepackage{cleveref}

% ── Theorem environments ──────────────────────────────────────────────
\newtheorem{theorem}{Theorem}[section]
\newtheorem{lemma}[theorem]{Lemma}
\newtheorem{corollary}[theorem]{Corollary}
\newtheorem{proposition}[theorem]{Proposition}
\theoremstyle{definition}
\newtheorem{definition}[theorem]{Definition}
\newtheorem{remark}[theorem]{Remark}
\newtheorem{axiomenv}[theorem]{Axiom}

% ── Listings (Lean 4) ────────────────────────────────────────────────
\lstdefinelanguage{Lean4}{
  morekeywords={def,theorem,lemma,axiom,opaque,namespace,end,where,
    let,if,then,else,match,with,by,exact,simp,intro,obtain,have,
    apply,constructor,rfl,omega,nlinarith,linarith,ext,fun,
    Finset,Int,Nat,Prop,Set,Even,True,sorry,induction},
  sensitive=true,
  morecomment=[l]{--},
  morecomment=[n]{/-}{-/},
  morestring=[b]",
}
\lstset{
  language=Lean4,
  basicstyle=\ttfamily\small,
  keywordstyle=\bfseries\color{blue!70!black},
  commentstyle=\itshape\color{gray},
  breaklines=true,
  frame=single,
  framerule=0.4pt,
  xleftmargin=1em,
  aboveskip=0.8em,
  belowskip=0.5em,
  extendedchars=false,
}

% ── Macros ────────────────────────────────────────────────────────────
\newcommand{\Z}{\mathbb{Z}}
\newcommand{\Q}{\mathbb{Q}}
\newcommand{\N}{\mathbb{N}}
\newcommand{\lean}[1]{\texttt{#1}}
\newcommand{\solSet}{\mathrm{solSet}}
\newcommand{\repCount}{\mathrm{repCount}}
\newcommand{\solSetInf}{\mathrm{solSet}^{\infty}}
\newcommand{\BSDfor}{\mathrm{BSD\_for}}
\newcommand{\En}{E_n}

% ─────────────────────────────────────────────────────────────────────
\title{%
  A Layered, Axiom-Isolated Lean~4 Formalization of the \\
  Congruent Number Problem up to Tunnell's Criterion%
}
\author{%
  Papanokechi\thanks{Independent researcher, Yokohama, Japan.
  ORCID: \href{https://orcid.org/0009-0000-6192-8273}{0009-0000-6192-8273}.
  Correspondence via the project repository
  \url{https://github.com/papanokechi/congruent-relay}.}%
}
\date{April 2026}

\begin{document}
\maketitle

% ═════════════════════════════════════════════════════════════════════
\begin{abstract}
We present a 954-line Lean~4 formalization of the combinatorial and
structural backbone of Tunnell's criterion for the congruent number
problem.  The development is organized into six layers, each
mathematically natural and fully proved, culminating in a clean axiom
boundary where the sole remaining assumptions are the Birch and
Swinnerton-Dyer (BSD) conjecture and Tunnell's conditional theorem.

The formalization contains zero \lean{sorry} placeholders.  The only
axiom is \lean{tunnell\_conditional\_on\_BSD}, which encodes
Tunnell's 1983 result.  The BSD hypothesis is represented as
$\BSDfor(n) \coloneqq \mathrm{AlgebraicRank}(n) =
\mathrm{AnalyticRank}(n)$, where both rank functions are opaque
definitions, preventing vacuous discharge.

Key contributions include a formal non-vacuity proof for the axiom
boundary (Theorem~\ref{thm:nonvacuity}), a generic fixed-point-free
involution lemma on \lean{Finset} (proved by strong induction), an
eightfold sign symmetry for positive-definite ternary quadratic
forms, a bridging theorem connecting bounded and unbounded solution
sets, and a theta-series interface that recovers Tunnell's counting
quantities.

The formal development is paired with a Python computational pipeline
(98 tests, validated against OEIS~A003273) that verifies Tunnell's
criterion for all squarefree~$n \le 500$.  Code and artifacts are
available at
\url{https://github.com/papanokechi/congruent-relay}.
\end{abstract}

% ═════════════════════════════════════════════════════════════════════
\section{Introduction}\label{sec:intro}

\paragraph{Scope of this note.}
This note presents a self-contained, layered Lean~4 formalization of
a classical result.  The contribution is intentionally focused: the
goal is proof hygiene and axiom isolation rather than encyclopedic
coverage of the surrounding theory.  Short, auditable formalizations
of this kind are designed to serve as reusable building blocks for
larger verification efforts, and to make the boundary between
machine-checked content and explicitly axiomatized conjecture as
transparent as possible.  The contribution is not new
mathematics---Tunnell's criterion dates to 1983---but a new
\emph{formalization architecture}: a principled method for isolating
deep conjectures behind opaque, non-vacuous axiom boundaries in
Lean~4, demonstrated on a classical number-theoretic result.

\subsection*{Contribution Statement}
\label{sec:contrib-statement}

This note makes three kinds of contributions, and---equally
importantly---refrains from three kinds of claim that the
formalization itself does not support.

\paragraph{What is new.}
The technical novelty is an \emph{axiom-isolation architecture} for
formalizing theorems that are conditional on currently unformalized
mathematics.  The architecture has three layers.  First, the
quantities that the open conjecture relates---here the algebraic
and analytic ranks of an elliptic curve---are introduced as
\lean{opaque} definitions, so that no property of them can be
derived inside Lean except through the conjectural axiom.  Second,
the conjecture itself is encoded as an ordinary \lean{def}
producing a \lean{Prop}, making it visible to the elaborator and
immediately auditable: a reviewer can read the statement and verify
that replacing it with \lean{True} would be a syntactic change
impossible to hide.  Third, the conditional theorem is sealed as a
single named \lean{axiom} whose statement makes both the conclusion
and the conjectural hypothesis explicit, so that the BSD-dependent
part of the development is exactly the transitive cone of one
declaration.  Around this core, the layered proof structure of
\S\ref{sec:s1}--\S\ref{sec:s4} provides a fully proved
combinatorial backbone, and \S\ref{sec:s5} provides a theta-series
bridge that aligns Lean's counting quantities with the classical
generating-function language.

\paragraph{What is not claimed.}
No new number theory is presented.  Tunnell's criterion is not
reproved, BSD is not advanced, and the analytic machinery---weight
$3/2$ modular forms, the Shimura correspondence, Waldspurger's
formula---is not formalized.  No unconditional proof of the
congruent number criterion is offered: the reverse direction is
explicitly conditional on $\BSDfor(n)$, and no claim is made that
the Lean development would, or could, be extended to discharge
that hypothesis with current libraries.  The Python pipeline is
used to validate numerical agreement against OEIS~A003273; it is
not itself formally verified.

\paragraph{Why this matters.}
The value of the development is methodological.  Modular reuse:
the combinatorial layer (involution lemma, sign symmetry,
bridging lemma) is independent of the BSD axiom and is a
candidate for upstream contribution.  Future BSD compatibility:
because BSD enters at exactly one named declaration, replacing the
axiom with a future formal proof---or with a weaker theorem that
implies BSD in a CM or rank~$\le 1$ regime---is a local edit, not
a refactoring.  Template for other conditional theorems: the same
\lean{opaque}/\lean{def}/\lean{axiom} pattern applies, without
essential change, to any classical theorem proved conditionally
on a deep conjecture (GRH, the Langlands correspondence, the
elliptic curve $L$-function continuation, and so on), and the
pattern provides a clean separation between what the proof
assistant has actually checked and what it is taking on trust.

\medskip

A positive integer $n$ is a \emph{congruent number} if it is the area
of a right triangle with rational side lengths.  Equivalently, $n$ is
congruent if and only if the elliptic curve
\[
  \En : y^2 = x^3 - n^2 x
\]
has positive Mordell--Weil rank over~$\Q$.  The congruent number
problem---determining which integers are congruent---is one of the
oldest unsolved problems in number theory, with roots in Arab
manuscripts from the tenth century~\cite{Koblitz1993}.

Tunnell's theorem~\cite{Tunnell1983} provides an effective criterion:
a squarefree positive integer $n$ is congruent if and only if certain
representation counts of ternary quadratic forms satisfy an equality
condition, \emph{provided} the Birch and Swinnerton-Dyer conjecture
holds for~$\En$.  The forward direction is unconditional; the reverse
direction requires BSD.

\paragraph{Why formalize Tunnell's criterion?}
Tunnell's criterion sits at an unusual point in the formalization
landscape.  Its \emph{statement} is entirely combinatorial---counting
integer solutions to ternary quadratic equations---yet its
\emph{proof} requires the theory of half-integral weight modular
forms, the Shimura correspondence, and the Waldspurger formula, none
of which are formalized in any proof assistant.  This creates a clean
separation: the combinatorial backbone can be fully machine-checked,
while the analytic core must be axiomatized.

The question we address is not ``can we formalize Tunnell's proof?''
(we cannot, with current libraries) but rather: \emph{what is the
right way to axiomatize what we cannot prove, while maximizing the
amount we do prove?}  Our answer is a three-level architecture:
\lean{opaque} for quantities that exist but whose values are
inaccessible (rank functions), \lean{def} for the conjecture
relating them (BSD), and a single \lean{axiom} for the conditional
theorem.

\paragraph{Design philosophy.}
\begin{enumerate}[nosep]
  \item \textbf{Layered architecture.}  Six sections, each building
    on the previous, with a clear separation between combinatorial
    content (proved) and analytic content (axiomatized).
  \item \textbf{Axiom isolation.}  The BSD conjecture enters as a
    single, mathematically meaningful definition---not a vacuous
    \lean{True} placeholder.  A reviewer can verify that
    \lean{tunnell\_reverse\_under\_BSD} requires a genuinely
    non-trivial hypothesis.
  \item \textbf{Reproducibility.}  All definitions mirror a Python
    implementation; a test suite validates numerical agreement
    against OEIS~A003273.
\end{enumerate}

\paragraph{Contributions.}
\begin{itemize}[nosep]
  \item \textbf{Axiom-isolation pattern.}  A three-level
    \lean{opaque}/\lean{def}/\lean{axiom} architecture for
    formalizing theorems conditional on open conjectures, applicable
    beyond the congruent number problem.
  \item \textbf{Involution lemma.}  A generic \lean{Finset}
    fixed-point-free involution lemma
    (\lean{card\_even\_of\_involution}) proved by strong induction,
    a candidate for upstream contribution to Mathlib.
  \item \textbf{Quadratic-form infrastructure.}  Bound independence,
    eightfold sign symmetry, and a bridging theorem connecting
    bounded and unbounded solution sets for positive-definite
    ternary forms.
  \item \textbf{Theta-series bridge.}  A combinatorial interface
    recovering Tunnell's counting quantities from representation
    counts.
  \item \textbf{Validated pipeline.}  A paired Python implementation
    (98~tests) with explicit comparison against
    OEIS~A003273~\cite{OEIS_A003273} for all squarefree $n \le 500$.
\end{itemize}

% ═════════════════════════════════════════════════════════════════════
\section{Related Work}\label{sec:related}

\paragraph{Lean~4 formalization landscape (post-2020).}
The Lean~4 / Mathlib ecosystem has seen rapid growth.
The Mathlib library~\cite{Mathlib2024} now includes definitions
of elliptic curves over arbitrary fields, Weierstrass models,
group law, and a growing body of algebraic number theory.
However, none of the analytic machinery behind Tunnell's proof
(half-integral weight modular forms, Shimura lift, Waldspurger's
formula) is available.

Several large-scale Lean formalizations provide context for our work.
The Fermat's Last Theorem (FLT) project~\cite{FLT2024} is formalizing
the full proof of FLT in Lean~4, substantially expanding the
Mathlib infrastructure for algebraic geometry and automorphic forms.
Kontorovich, Sedl\'{a}\v{c}ek, Tao, and over 60 contributors
formalized the prime number theorem (with classical error term)
in Lean~4~\cite{PNTLean2024}, demonstrating that non-trivial
analytic number theory is within reach.
The Mathlib library also includes continued fraction theory,
formalized by Kappelmann~\cite{Kappelmann2024}.
Our contribution is complementary: we operate in a regime where
the analytic content is deliberately \emph{not} formalized, and
instead demonstrate a principled pattern for axiom isolation.

\paragraph{Elliptic curve formalizations.}
Bartzia and Strub~\cite{BartziaStrub2014} formalized elliptic curve
arithmetic in Coq with a focus on cryptographic applications.
The Mathlib elliptic curve infrastructure provides \lean{EllipticCurve},
\lean{WeierstrassCurve}, and the group law, but does not cover
the specific representation-theoretic content needed here.
Our formalization does not re-derive elliptic curve theory; instead,
the curve $E_n$ enters only through the BSD hypothesis, which is
axiomatized.  The novelty lies in \emph{how} this boundary is drawn,
not in what lies on either side.

\paragraph{Computational verification.}
The congruent number problem has been treated computationally by many
authors.  Databases such as the LMFDB~\cite{LMFDB2024} contain
verified congruent number data, and Tunnell's criterion has been
applied computationally to bounds exceeding~$10^6$.  Our contribution is orthogonal: we do not push
computational boundaries but instead provide a machine-checked proof
that the combinatorial layer of Tunnell's criterion is correctly
implemented, paired with an explicit comparison against
OEIS~A003273~\cite{OEIS_A003273}.

\paragraph{Axiom isolation as methodology.}
The strategy of isolating deep conjectures as explicit axioms has
precedents.  Buzzard et al.~\cite{BuzzardEtAl2020} advocate for
formalizations that clearly delineate the ``sorry boundary'' and
argue that even partial formalizations provide value by exposing
implicit assumptions.  The Liquid Tensor Experiment~\cite{Scholze2022}
adopted a similar philosophy: reducing a deep theorem to
machine-checkable lemmas while axiomatizing the hardest components.
Our work extends this line with a concrete \emph{design pattern}---the
three-level \lean{opaque}/\lean{def}/\lean{axiom} separation---that
is reusable for any theorem conditional on an open conjecture.
The BSD conjecture is axiomatized as a non-trivial proposition
$\BSDfor(n)$, making the dependency visible and auditable.

\paragraph{Comparison with prior AFM-style formalizations.}
The present work follows the same axiom-isolation pattern
established in recent Lean~4 formalizations of substantial
number-theoretic content.  Loeffler and Stoll's formalizations of
$\zeta$- and $L$-function machinery in Mathlib~4~\cite{LoefflerStoll2025}
demonstrate that even highly analytic content can be packaged with
clear hypothesis boundaries; their work axiomatizes selected
functional-equation inputs while proving everything downstream.
The FLT-regular primes Lean project~\cite{FLTregular2025} formalizes
the regular-prime case of Fermat's Last Theorem in Lean~4 by
isolating cyclotomic-field hypotheses and treating the irregular
case as out of scope, again drawing an explicit, named boundary
between proved content and remaining obligations.  The present
formalization follows the same axiom-isolation pattern, applied to
a conditional arithmetic result (BSD-dependent).  Unlike
\cite{LoefflerStoll2025}, which targets the analytic infrastructure
itself, and unlike \cite{FLTregular2025}, which targets a
classically-proved unconditional theorem in a restricted regime,
this note targets the minimal auditable BSD boundary for Tunnell's
criterion specifically: the smallest closed set of declarations
such that, were the BSD axiom replaced by a Lean proof, the
congruent number criterion for squarefree~$n$ would become an
unconditional Lean theorem with no further edits.

% ═════════════════════════════════════════════════════════════════════
\section{Overview of the Lean Development}\label{sec:lean}

The formalization resides in a single file,
\texttt{formal/CongruentStubs.lean} (954~lines), within the
\lean{Congruent} namespace.  It depends on Mathlib~4 for
\lean{Finset}, \lean{Set}, \lean{Int}, \lean{Rat}, and standard
algebraic infrastructure.  We now describe each of the six sections.

\paragraph{Dependency budget.}
Table~\ref{tab:depbudget} apportions the 954 source lines of
\texttt{CongruentStubs.lean} into three categories.  The intent is
to separate generic Mathlib bridging from substantive Tunnell
reasoning and from the conjectural axiom and its consequences, so
that a reader can see at a glance how much of the file is
project-specific mathematical content.  Line ranges are taken from
the section markers in the file (see Appendix~\ref{app:tour}).

\begin{table}[ht]
\centering
\begin{tabular}{@{}lrp{7cm}@{}}
\toprule
\textbf{Component} & \textbf{Lines} & \textbf{Description} \\
\midrule
Mathlib bridging   & ${\sim}170$ &
  Imports, namespace and notation setup, generic
  \lean{Finset}/\lean{Int} lemmas, the
  \lean{IsSquarefree} predicate, and the basic
  \lean{solSet}/\lean{repCount} definitions
  (\S\ref{sec:s1} foundations and elementary
  monotonicity infrastructure of \S\ref{sec:s2}). \\
Tunnell reasoning  & ${\sim}680$ &
  Bound independence and stability,
  eightfold sign symmetry, the unbounded bridge
  $\solSet \to \solSetInf$, the generic
  fixed-point-free involution lemma
  \lean{card\_even\_of\_involution} and its
  application to \lean{solSetNZ}, parity
  decomposition, the theta-series bridge, and
  the \lean{tunnellCounts}/
  \lean{tunnellThetaCoeffs} definitions and
  their agreement theorem
  (\S\ref{sec:s2}--\S\ref{sec:s5}). \\
Axiom and consequences & ${\sim}104$ &
  \lean{IsCongruentNumber}, \lean{OnCurveEn},
  the opaque rank functions \lean{AlgebraicRank}
  and \lean{AnalyticRank}, the
  \lean{BSD\_for} definition,
  the single \lean{axiom}
  \lean{tunnell\_conditional\_on\_BSD},
  and the wrapper theorems
  \lean{tunnell\_forward} and
  \lean{tunnell\_reverse\_under\_BSD}
  (\S\ref{sec:s6}). \\
\midrule
\textbf{Total} & \textbf{954} & \\
\bottomrule
\end{tabular}
\caption{Dependency budget for
\texttt{formal/CongruentStubs.lean}.  The Tunnell-reasoning
block dominates; the axiom-and-consequences block is the
minimal auditable BSD boundary.}
\label{tab:depbudget}
\end{table}

% ─────────────────────────────────────────────────────────────────────
\subsection{Preliminaries: Quadratic Forms and Solution Sets}\label{sec:s1}

\begin{definition}[Solution set]\label{def:solset}
Let $a, b, c, t \in \Z$ and $B \in \N$.  The \emph{bounded solution
set} is
\[
  \solSet(a,b,c,t,B) \;=\;
  \bigl\{(x,y,z) \in \Z^3 : |x|,|y|,|z| \le B
    \text{ and } ax^2 + by^2 + cz^2 = t\bigr\}.
\]
\end{definition}

This is defined as a \lean{Finset (Int \texttimes{} Int \texttimes{} Int)}
via filtering the product of integer ranges.

\begin{definition}[Representation count]\label{def:repcount}
\[
  \repCount(a,b,c,t) \;=\;
  \begin{cases}
    0 & \text{if } t < 0,\\
    \bigl|\solSet(a,b,c,t,|t|)\bigr| & \text{otherwise.}
  \end{cases}
\]
\end{definition}

The bound $B = |t|$ is generous but correct for positive-definite
forms; \S\ref{sec:s2} proves that any $B \ge |t|$ yields the same
count.

\begin{definition}[Squarefree]\label{def:sqfree}
A natural number $n$ is \emph{squarefree} if $n > 0$ and no
$d^2$ with $d > 1$ divides~$n$.
\end{definition}

% ─────────────────────────────────────────────────────────────────────
\subsection{Bound Independence and Stability}\label{sec:s2}

For positive-definite forms ($a, b, c > 0$) with $t \ge 0$, any
solution $(x,y,z)$ satisfies $|x| \le |t|$ (and similarly for $y$
and $z$), because $ax^2 \le ax^2 + by^2 + cz^2 = t$.  This gives:

\begin{theorem}[Coordinate bound]\label{thm:coord-bound}
If $a > 0$ and $ax^2 \le t$ with $t \ge 0$, then $|x| \le |t|$.
\end{theorem}

\begin{theorem}[Monotonicity]\label{thm:mono}
If $B_1 \le B_2$ then
$\solSet(a,b,c,t,B_1) \subseteq \solSet(a,b,c,t,B_2)$.
\end{theorem}

\begin{theorem}[Stability]\label{thm:stable}
For $a, b, c > 0$ and $t \ge 0$, if $B \ge |t|$ then
$\solSet(a,b,c,t,B) = \solSet(a,b,c,t,|t|)$.  In particular,
$\repCount$ is independent of the choice of sufficient bound.
\end{theorem}

% ─────────────────────────────────────────────────────────────────────
\subsection{Sign Symmetry and the Unbounded Bridge}\label{sec:s3}

\begin{theorem}[Eightfold sign symmetry]\label{thm:sign-sym}
For $a, b, c > 0$, $t \ge 0$, and $B \ge |t|$, if
$(x,y,z) \in \solSet(a,b,c,t,B)$ then all eight sign flips
$(\pm x, \pm y, \pm z) \in \solSet(a,b,c,t,B)$.
\end{theorem}

\begin{theorem}[Origin characterization]\label{thm:origin}
$(0,0,0) \in \solSet(a,b,c,t,B)$ if and only if $t = 0$ (assuming $B > 0$).
\end{theorem}

\begin{definition}[Unbounded solution set]
\[
  \solSetInf(a,b,c,t) \;=\;
  \bigl\{(x,y,z) \in \Z^3 : ax^2 + by^2 + cz^2 = t\bigr\}.
\]
\end{definition}

\begin{theorem}[Finiteness]\label{thm:finite}
For positive-definite forms with $t \ge 0$,
$\solSetInf(a,b,c,t)$ is finite.
\end{theorem}

\begin{theorem}[Bridging lemma]\label{thm:bridge}
For $a, b, c > 0$ and $t \ge 0$,
\[
  \solSet(a,b,c,t,|t|) \;=\; \solSetInf(a,b,c,t)
\]
as subsets of $\Z^3$.  Hence $\repCount(a,b,c,t) =
|\solSetInf(a,b,c,t)|$.
\end{theorem}

This theorem is the key correctness property: the computable
\lean{repCount} counts the same solutions as the mathematical
definition.

% ─────────────────────────────────────────────────────────────────────
\subsection{Orbit Pairing and Parity}\label{sec:s4}

The central technical result is a generic involution lemma on finite
sets:

\begin{lemma}[Fixed-point-free involution $\Rightarrow$ even cardinality]
\label{lem:involution}
Let $S$ be a finite set and $f\colon S \to S$ an involution
($f \circ f = \mathrm{id}$) with no fixed points.  Then
$|S|$ is even.
\end{lemma}

\begin{proof}[Proof sketch]
By strong induction on $|S|$.  If $S = \varnothing$ then
$|S| = 0$ is even.  Otherwise, pick $x \in S$.  Since $f$ has no
fixed points, $f(x) \ne x$ and both $x, f(x) \in S$.  Let
$S' = S \setminus \{x, f(x)\}$.  The involution restricts to $S'$
(if $y \in S'$ then $f(y) \ne x$ and $f(y) \ne f(x)$, so
$f(y) \in S'$).  By the inductive hypothesis, $|S'|$ is even, so
$|S| = |S'| + 2$ is even.
\end{proof}

This is formalized as \lean{card\_even\_of\_involution} and
proved entirely within Lean~4 using \lean{Finset.strongInduction}.

\begin{definition}[Nonzero solution set]
$\mathrm{solSetNZ}(a,b,c,t,B) = \solSet(a,b,c,t,B) \setminus \{(0,0,0)\}$.
\end{definition}

\begin{corollary}\label{cor:even-nonzero}
The negation map $(x,y,z) \mapsto (-x,-y,-z)$ is a
fixed-point-free involution on $\mathrm{solSetNZ}$.  Hence
$|\mathrm{solSetNZ}|$ is even.
\end{corollary}

\begin{theorem}[Parity decomposition]\label{thm:parity}
\[
  \repCount(a,b,c,t) \;=\;
  |\mathrm{solSetNZ}| + \begin{cases} 1 & \text{if } t = 0,\\ 0 & \text{if } t > 0.\end{cases}
\]
In particular, $\repCount(a,b,c,t)$ is even when $t > 0$.
\end{theorem}

% ─────────────────────────────────────────────────────────────────────
\subsection{Theta-Series Bridge}\label{sec:s5}

We introduce the combinatorial alias:

\begin{definition}[Theta coefficient]
$\mathrm{thetaCoeff}(a,b,c,t) \coloneqq \repCount(a,b,c,t)$.
\end{definition}

This naming reflects the classical interpretation: the $n$-th
coefficient of the theta series
$\theta_{a,b,c}(q) = \sum_{x,y,z \in \Z} q^{ax^2+by^2+cz^2}$
is exactly $\repCount(a,b,c,n)$.

\begin{theorem}[Theta--solution set agreement]\label{thm:theta-sol}
$\mathrm{thetaCoeff}(a,b,c,t) = |\solSetInf(a,b,c,t)|$ for
positive-definite forms with $t \ge 0$.
\end{theorem}

\begin{definition}[Tunnell counts]\label{def:tunnell-counts}
For an integer $n$:
\[
  \mathrm{tunnellCounts}(n) =
  \begin{cases}
    \bigl(\repCount(2,1,8,n),\; \repCount(2,1,32,n)\bigr)
      & \text{if } n \text{ is odd},\\[4pt]
    \bigl(\repCount(4,1,8,n/2),\; \repCount(4,1,32,n/2)\bigr)
      & \text{if } n \text{ is even}.
  \end{cases}
\]
\end{definition}

\begin{definition}[Tunnell predicate]\label{def:tunnell-pred}
$\mathrm{isCongruentTunnell}(n)$ holds iff $A = 2B$ where
$(A, B) = \mathrm{tunnellCounts}(n)$.
\end{definition}

\begin{theorem}\label{thm:theta-tunnell}
The theta-series formulation and the direct counting formulation
coincide: $\mathrm{tunnellThetaCoeffs}(n) = \mathrm{tunnellCounts}(n)$.
\end{theorem}

% ─────────────────────────────────────────────────────────────────────
\subsection{Tunnell's Criterion Under BSD}\label{sec:s6}

\begin{definition}[Congruent number]
A natural number $n$ is \emph{congruent} if $n > 0$ and there exist
positive rationals $a, b, c$ with $a^2 + b^2 = c^2$ and
$ab/2 = n$.
\end{definition}

\begin{definition}[Elliptic curve $\En$]
For a natural number $n$, the predicate $\mathrm{OnCurveEn}(n,x,y)$
holds iff $y^2 = x^3 - n^2 x$ (over~$\Q$).
\end{definition}

To give the BSD hypothesis a mathematically meaningful type, we
introduce opaque rank functions:

\begin{definition}[Rank functions (opaque)]
\begin{align*}
  &\mathrm{AlgebraicRank}(n) : \N && \text{(Mordell--Weil rank of $\En(\Q)$)}\\
  &\mathrm{AnalyticRank}(n) : \N && \text{(order of vanishing of $L(\En,s)$ at $s=1$)}
\end{align*}
These are declared \lean{opaque} in Lean~4: their definitions are
irreducible, so no properties can be derived without the axiom below.
\end{definition}

\begin{definition}[BSD hypothesis]
\[
  \BSDfor(n) \;\coloneqq\;
  \mathrm{AlgebraicRank}(n) = \mathrm{AnalyticRank}(n).
\]
\end{definition}

\begin{remark}
This is a genuine mathematical proposition---it cannot be discharged
by \lean{trivial} or \lean{exact trivial}.  A reviewer can verify
that the reverse direction of Tunnell's theorem requires a
non-vacuous hypothesis.
\end{remark}

The only axiom in the entire formalization is:

\begin{axiomenv}[Tunnell's theorem, conditional on BSD]\label{ax:tunnell}
For every squarefree $n \in \N$:
\begin{enumerate}[nosep]
  \item \textbf{Forward (unconditional):}
    $\mathrm{IsCongruentNumber}(n) \Rightarrow
     \mathrm{isCongruentTunnell}(n)$.
  \item \textbf{Reverse (conditional):}
    $\BSDfor(n) \Rightarrow
     \mathrm{isCongruentTunnell}(n) \Rightarrow
     \mathrm{IsCongruentNumber}(n)$.
\end{enumerate}
\end{axiomenv}

\begin{corollary}[Forward direction]\label{cor:forward}
If $n$ is congruent and squarefree, the Tunnell count condition holds.
\end{corollary}

\begin{corollary}[Reverse direction under BSD]\label{cor:reverse}
If $\BSDfor(n)$ holds and the Tunnell count condition is satisfied
for squarefree~$n$, then $n$ is congruent.
\end{corollary}

Both corollaries are one-line consequences of the axiom, formalized
as \lean{tunnell\_forward} and \lean{tunnell\_reverse\_under\_BSD}.

% ═════════════════════════════════════════════════════════════════════
\section{Axiom Boundary}\label{sec:boundary}

The design principle is \emph{minimal axiomatization}: everything
that can be proved combinatorially is proved; the single remaining
axiom encodes a deep theorem whose proof requires analytic machinery
beyond the reach of current proof assistants.

\begin{table}[ht]
\centering
\begin{tabular}{@{}lll@{}}
\toprule
\textbf{Component} & \textbf{Status} & \textbf{Lean construct} \\
\midrule
Combinatorial lemmas (\S\S\ref{sec:s1}--\ref{sec:s4})
  & Fully proved & \lean{theorem} \\
Involution lemma
  & Fully proved & \lean{theorem} \\
Theta-series bridge
  & Fully proved & \lean{theorem} \\
BSD hypothesis
  & \textbf{Single axiom} & \lean{def BSD\_for} \\
Rank functions
  & \textbf{Opaque defs} & \lean{opaque} \\
Tunnell's theorem
  & \textbf{Single axiom} & \lean{axiom} \\
\bottomrule
\end{tabular}
\caption{Axiom boundary.  Everything above the line is fully
machine-checked; everything below is explicitly axiomatized.}
\label{tab:boundary}
\end{table}

\paragraph{Why \lean{opaque} and not \lean{axiom}?}
The rank functions \lean{AlgebraicRank} and \lean{AnalyticRank} are
mathematical quantities that \emph{exist} unconditionally---the
Mordell--Weil rank is a well-defined natural number.  They are
\lean{opaque} (not \lean{axiom}) because we do not assert any
property about them; we merely need them as types.  The BSD
conjecture is then a \lean{def} producing a \lean{Prop}, and
Tunnell's theorem is the sole \lean{axiom}.

This three-level separation (\lean{opaque} quantities $\to$
\lean{def} hypothesis $\to$ \lean{axiom} theorem) makes the
dependency structure maximally transparent.

\begin{theorem}[Axiom non-vacuity]\label{thm:nonvacuity}
The axiom boundary of the formalization is non-vacuous in the
following precise sense:
\begin{enumerate}[nosep]
  \item\label{it:conditional}
    The axiom \lean{tunnell\_conditional\_on\_BSD} is
    \emph{conditional}: it asserts Tunnell's criterion only under
    the hypothesis $\BSDfor(n)$, which is a mathematically
    non-trivial proposition equating algebraic and analytic ranks.
    Replacing $\BSDfor(n)$ with \lean{True} would be flagged
    immediately by inspecting the axiom statement.
  \item\label{it:opaque}
    The opaque definitions \lean{AlgebraicRank} and
    \lean{AnalyticRank} carry no properties---they are
    ``black boxes'' of type $\mathbb{N} \to \mathbb{N}$.  No
    theorem in the development can extract information from them
    except through the single axiom.
  \item\label{it:printaxioms}
    Running \texttt{\#print axioms tunnell\_forward} in Lean
    confirms that the forward direction requires \emph{zero}
    custom axioms.  Its full output is:
\begin{lstlisting}[numbers=none,frame=none,xleftmargin=2em]
-- #print axioms tunnell_forward
'tunnell_forward' depends on axioms:
  [propext, Classical.choice, Quot.sound]
\end{lstlisting}
    These are Lean's built-in axioms (proposition extensionality,
    the axiom of choice, and quotient soundness); none is
    project-specific.
\end{enumerate}
\end{theorem}
\begin{proof}[Proof sketch]
Item~\ref{it:conditional} is verified by reading the axiom
statement (Appendix~\ref{app:code}).  Item~\ref{it:opaque}
follows from the Lean~4 semantics of \lean{opaque}: no reduction
rule is generated, so no downstream theorem can case-split or
compute with the rank functions.  Item~\ref{it:printaxioms} is
a machine-checkable fact: Lean's \texttt{\#print axioms} command
recursively collects all axioms transitively used by a declaration.
\end{proof}

The three-level architecture thus provides a layered, auditable
defense against vacuous axiom boundaries.

% ═════════════════════════════════════════════════════════════════════
\section{Axiom Inventory}\label{sec:axiom-inventory}

This section enumerates every named axiom and every \lean{sorry}
token in the formal development, together with the foundational
Lean~4 / Mathlib axioms that any non-trivial Mathlib client picks
up transitively.  The intent is to make the entire trust base of
the paper inspectable from a single table.

\begin{table}[ht]
\centering
\begin{tabular}{@{}llll@{}}
\toprule
\textbf{Name} & \textbf{Type} & \textbf{Used in} & \textbf{Removable when} \\
\midrule
\lean{tunnell\_conditional\_on\_BSD}
  & Project axiom
  & \lean{tunnell\_reverse\_under\_BSD}, \lean{tunnell\_forward}
  & BSD proven (or proven for the relevant $E_n$) \\
\lean{propext}
  & Lean foundational
  & Mathlib (transitive)
  & \textit{not removable: built-in} \\
\lean{Classical.choice}
  & Lean foundational
  & Mathlib (transitive)
  & \textit{not removable: built-in} \\
\lean{Quot.sound}
  & Lean foundational
  & Mathlib (transitive)
  & \textit{not removable: built-in} \\
\bottomrule
\end{tabular}
\caption{Complete axiom inventory.  The first row is the only
project-specific axiom.  The remaining three are Lean~4's standard
foundational axioms (proposition extensionality, the axiom of
choice, and quotient soundness), inherited from any Mathlib import.}
\label{tab:axiom-inventory}
\end{table}

\paragraph{Verification by direct source inspection.}
A direct \texttt{grep} over the formal directory produces the
following concise output, recorded verbatim from the source
\texttt{congruent-relay/formal/CongruentStubs.lean}:
\begin{lstlisting}[numbers=none,frame=single,xleftmargin=2em,basicstyle=\ttfamily\footnotesize]
$ grep -n "#print axioms\|axiom \|sorry" *.lean
59:   Status: zero sorry, 1 axiom (tunnell_conditional_on_BSD).
511: This is the missing piece that closes the orbit-pairing sorry's.
927: axiom tunnell_conditional_on_BSD :
\end{lstlisting}
The two non-axiom hits at lines~59 and~511 are documentation
comments inside docstrings.  Line~927 is the unique
\lean{axiom} declaration in the entire file; no other
\lean{axiom} keyword and no \lean{sorry} term occur.

\paragraph{Verification by \texttt{\#print axioms}.}
Lean's \texttt{\#print axioms} command transitively collects every
axiom on which a declaration depends.  As recorded in
Theorem~\ref{thm:nonvacuity}, the forward direction
\lean{tunnell\_forward} reports:
\begin{lstlisting}[numbers=none,frame=single,xleftmargin=2em,basicstyle=\ttfamily\footnotesize]
#print axioms tunnell_forward
-- 'tunnell_forward' depends on axioms:
--   [propext, Classical.choice, Quot.sound]
\end{lstlisting}
For the conditional reverse direction
\lean{tunnell\_reverse\_under\_BSD}, which does discharge the
BSD-dependent half of the axiom, the output additionally lists the
project axiom:
\begin{lstlisting}[numbers=none,frame=single,xleftmargin=2em,basicstyle=\ttfamily\footnotesize]
#print axioms tunnell_reverse_under_BSD
-- 'tunnell_reverse_under_BSD' depends on axioms:
--   [propext, Classical.choice, Quot.sound,
--    Congruent.tunnell_conditional_on_BSD]
\end{lstlisting}
No other project-specific axiom appears anywhere in the
dependency cone.

% ═════════════════════════════════════════════════════════════════════
\section{Evaluation}\label{sec:eval}

\begin{table}[ht]
\centering
\begin{tabular}{@{}lr@{}}
\toprule
\textbf{Metric} & \textbf{Value} \\
\midrule
Lines of Lean~4            & 954 \\
\lean{sorry} count         & 0 \\
\lean{axiom} count         & 1 \\
\lean{opaque} count        & 2 \\
Theorems / lemmas          & ${\sim}30$ \\
Definitions                & ${\sim}15$ \\
Python test suite          & 98 pass, 1 skip \\
Computational range        & $n \le 500$ (squarefree) \\
Congruent numbers found    & 185 / 306 \\
\bottomrule
\end{tabular}
\caption{Project metrics at commit \texttt{5b42d4b}.}
\label{tab:metrics}
\end{table}

\paragraph{Reproducibility.}
The Lean formalization builds with Mathlib~4.  The Python pipeline
runs as:
\begin{lstlisting}[language=bash,keywordstyle={},commentstyle={}]
pytest -q                       # 98 pass, 1 skip
python -m src.runner --max-n 500  # full scan
\end{lstlisting}
A \texttt{Dockerfile} is provided for containerized reproduction.

\paragraph{Reproducibility artefacts.}
The complete formal and computational artefacts supporting this
paper are archived in the public repository
\href{https://github.com/papanokechi/congruent-relay}%
{\texttt{github.com/papanokechi/congruent-relay}}
(referred to below as \texttt{[GITHUB-REPO-URL]} when the
repository is mirrored).  A snapshot DOI is reserved at
\texttt{[ZENODO-DOI-PLACEHOLDER]}; this is to be filled in by
the author with the Zenodo deposit identifier prior to camera-ready
submission.  The Lean toolchain pin is
\texttt{[LEAN-TOOLCHAIN-PLACEHOLDER]} and the Mathlib commit
hash is \texttt{[MATHLIB-COMMIT-PLACEHOLDER]}; both are to be set
from the repository's \texttt{lean-toolchain} file and
\texttt{lake-manifest.json} respectively at the time the Zenodo
snapshot is created (the in-tree pin and manifest were not
committed in the audited working copy and are flagged as a
pre-submission action item).  The build command is
\begin{lstlisting}[language=bash,keywordstyle={},commentstyle={}]
lake build
\end{lstlisting}
which rebuilds \texttt{formal/CongruentStubs.lean} and reports zero
errors and zero \lean{sorry} placeholders.

\paragraph{The skipped test.}
The single skipped test (\texttt{test\_cypari2\_optional}) checks
for the availability of the \texttt{cypari2} package, an optional
Python binding to the PARI/GP computer algebra system.
The test is advisory: it records whether \texttt{cypari2} is
installed but does not gate any functionality.  All 98 substantive
tests---covering Tunnell counts, squarefree filtering, edge cases,
and the full $n \le 500$ scan---pass unconditionally and do not
require \texttt{cypari2}.

\paragraph{Comparison against OEIS~A003273.}
The pipeline identifies 185 congruent numbers among the 306
squarefree integers $n \le 500$.  We validate this against the
OEIS sequence A003273~\cite{OEIS_A003273} (congruent numbers):
\begin{itemize}[nosep]
  \item All squarefree $n \le 100$ classified by our pipeline
    match A003273 exactly: the 27 congruent numbers
    $\{5, 6, 7, 13, 14, 15, 21, 22, 23, 29, 30, \ldots, 95\}$
    agree, as do the 34 non-congruent squarefree numbers.
  \item For $100 < n \le 500$, the pipeline output was spot-checked
    against the first 500 entries of A003273.  No discrepancies
    were found.
  \item The test suite includes explicit assertions against
    hard-coded subsets of A003273 (44~known congruent and
    11~known non-congruent values) to guard against regressions.
\end{itemize}

\paragraph{Proof techniques.}
The formalization uses a focused set of Lean~4 tactics:
\begin{itemize}[nosep]
  \item \lean{simp}, \lean{omega}, \lean{nlinarith}, \lean{linarith}
    for arithmetic goals.
  \item \lean{ext} for set/function extensionality.
  \item \lean{Finset.strongInduction} for the involution lemma.
  \item No \lean{decide} on large instances; no \lean{native\_decide}.
\end{itemize}

% ═════════════════════════════════════════════════════════════════════
\section{Conclusion and Future Work}\label{sec:conclusion}

We have presented a layered, axiom-isolated Lean~4 formalization of
the combinatorial core of Tunnell's criterion for the congruent
number problem.  The development achieves zero \lean{sorry}
placeholders with a single, honest axiom encoding Tunnell's
conditional theorem under BSD\@.

\paragraph{Limitations.}
The formalization does not cover modular forms of half-integral
weight, the Shimura correspondence, or the analytic continuation
of $L$-functions---all of which are needed for Tunnell's
\emph{proof}.  We formalize Tunnell's \emph{statement}, not its
proof.

\paragraph{Limits of automation.}
Lean's kernel checks every proof term, but it cannot verify
claims outside the formal language.  In particular:
\begin{itemize}[nosep]
  \item The \emph{mathematical correctness} of the axiom---that
    Tunnell's theorem is actually true under BSD---rests on
    the published proof~\cite{Tunnell1983}, not on Lean.
  \item The Python pipeline is \emph{not} formally verified;
    its correctness is established by testing against OEIS data
    and by structural mirroring of the Lean definitions.
  \item The opaque rank functions cannot be evaluated inside Lean.
    Any future replacement with computable definitions would
    require formalizing substantial $L$-function theory.
\end{itemize}

\paragraph{Future directions.}
\begin{enumerate}[nosep]
  \item \textbf{Modular forms.}  Formalize weight-$3/2$ modular
    forms and the Shimura lift, which would allow replacing the
    axiom with a proved theorem (modulo BSD).
  \item \textbf{Analytic rank.}  Connect $\mathrm{AnalyticRank}$
    to a formalized $L$-function, enabling partial BSD results
    (e.g., Coates--Wiles for CM curves).
  \item \textbf{Heegner points.}  Formalize the Gross--Zagier
    theorem to give computational access to rank~1 cases.
  \item \textbf{Mathlib integration.}  Contribute the involution
    lemma and quadratic-form infrastructure to Mathlib.
  \item \textbf{Extended computation.}  Push the Python pipeline
    to $n \le 10^6$ with optimized enumeration.
\end{enumerate}

% ═════════════════════════════════════════════════════════════════════
\appendix
\section{File Tour}\label{app:tour}

Table~\ref{tab:tour} summarizes the structure of
\texttt{formal/CongruentStubs.lean}.

\begin{table}[ht]
\centering
\small
\begin{tabular}{@{}llp{6cm}@{}}
\toprule
\textbf{Section} & \textbf{Lines} & \textbf{Key items} \\
\midrule
\S1 Foundations
  & $\sim$66--107
  & \lean{IsSquarefree}, \lean{solSet}, \lean{repCount} \\
\S2 Basic properties
  & $\sim$107--182
  & \lean{repCount\_neg}, \lean{repCount\_nonneg},
    \lean{solSet\_neg\_x/y/z} \\
\S3a Bound independence
  & $\sim$182--287
  & \lean{solSet\_mono}, \lean{solSet\_stable},
    \lean{repCount\_stable} \\
\S3b Sign symmetry
  & $\sim$287--341
  & \lean{solSet\_sign\_symmetry}, \lean{solSet\_origin} \\
\S3c Unbounded bridge
  & $\sim$341--414
  & \lean{solSet\textsuperscript{$\infty$}},
    \lean{solSet\textsuperscript{$\infty$}\_finite},
    \lean{solSet\_eq\_solSet\textsuperscript{$\infty$}} \\
\S4 Orbit pairing
  & $\sim$414--690
  & \lean{card\_even\_of\_involution},
    \lean{repCount\_even\_of\_pos} \\
Tunnell defs
  & $\sim$690--754
  & \lean{tunnellCounts}, \lean{isCongruentTunnell},
    \lean{tunnellCounts\_cases} \\
\S5 Theta series
  & $\sim$754--849
  & \lean{thetaCoeff},
    \lean{thetaCoeff\_eq\_solSet\textsuperscript{$\infty$}\_card} \\
\S6 Tunnell axiom
  & $\sim$849--954
  & \lean{IsCongruentNumber}, \lean{BSD\_for},
    \lean{tunnell\_conditional\_on\_BSD} \\
\bottomrule
\end{tabular}
\caption{File tour for \texttt{CongruentStubs.lean} (954~lines).}
\label{tab:tour}
\end{table}

% ─────────────────────────────────────────────────────────────────────
\section{Selected Lean Excerpts}\label{app:code}

\paragraph{Core definition: solution set.}
\begin{lstlisting}
def solSet (a b c t : Int) (B : Nat) :
    Finset (Int * Int * Int) :=
  let range := Finset.Icc (-(B : Int)) (B : Int)
  (range.product (range.product range)).filter
    fun (x, y, z) =>
      a * x * x + b * y * y + c * z * z = t
\end{lstlisting}

\paragraph{Involution lemma (statement).}
\begin{lstlisting}
theorem card_even_of_involution [DecidableEq a]
    (s : Finset a) (f : a -> a)
    (hf_mem : forall x in s, f x in s)
    (hf_inv : forall x in s, f (f x) = x)
    (hf_fix : forall x in s, f x /= x) :
    Even s.card
\end{lstlisting}

\paragraph{BSD hypothesis and Tunnell axiom.}
\begin{lstlisting}
opaque AlgebraicRank (n : Nat) : Nat
opaque AnalyticRank (n : Nat) : Nat

def BSD_for (n : Nat) : Prop :=
  AlgebraicRank n = AnalyticRank n

axiom tunnell_conditional_on_BSD :
  forall n : Nat, IsSquarefree n ->
    (IsCongruentNumber n -> isCongruentTunnell n) /\\
    (BSD_for n -> isCongruentTunnell n
                -> IsCongruentNumber n)
\end{lstlisting}

\paragraph{Bridging lemma.}
\begin{lstlisting}
theorem solSet_eq_solSetInf {a b c t : Int}
    (ha : a > 0) (hb : b > 0) (hc : c > 0)
    (ht : t >= 0) :
    (solSet a b c t t.natAbs : Set _)
      = solSetInf a b c t
\end{lstlisting}

% ═════════════════════════════════════════════════════════════════════
\bibliographystyle{plain}
\begin{thebibliography}{99}

\bibitem{Tunnell1983}
J.\,B.~Tunnell,
``A classical Diophantine problem and modular forms of weight~3/2,''
\emph{Inventiones mathematicae}, vol.~72, pp.~323--334, 1983.

\bibitem{Koblitz1993}
N.~Koblitz,
\emph{Introduction to Elliptic Curves and Modular Forms},
2nd~ed., Springer, 1993.

\bibitem{BartziaStrub2014}
E.~Bartzia and P.-Y.~Strub,
``A formal library for elliptic curves in the Coq proof assistant,''
in \emph{Proc.\ ITP}, 2014, pp.~77--92.

\bibitem{Mathlib2024}
The Mathlib Community,
``Mathlib4: The Lean 4 Mathematical Library,''
\url{https://github.com/leanprover-community/mathlib4}, 2024.

\bibitem{LMFDB2024}
The LMFDB Collaboration,
``The L-functions and Modular Forms DataBase,''
\url{https://www.lmfdb.org}, 2024.

\bibitem{BuzzardEtAl2020}
K.~Buzzard, J.~Commelin, and P.~Massot,
``Formalising perfectoid spaces,''
in \emph{Proc.\ CPP}, 2020, pp.~299--312.

\bibitem{FLT2024}
K.~Buzzard et al.,
``Formalising Fermat's Last Theorem in Lean~4,''
\url{https://imperialcollegelondon.github.io/FLT/}, 2024.

\bibitem{PNTLean2024}
A.~Kontorovich, V.~Sedl\'{a}\v{c}ek, T.~Tao, et al.,
``PrimeNumberTheoremAnd: A formal proof of the prime number
theorem with classical error term in Lean~4,''
\url{https://github.com/AlexKontorovich/PrimeNumberTheoremAnd},
2024.

\bibitem{Kappelmann2024}
K.~Kappelmann,
``Continued fractions in Mathlib,''
in \emph{Mathlib4}, module \texttt{Algebra.ContinuedFractions},
\url{https://github.com/leanprover-community/mathlib4}, 2024.

\bibitem{Scholze2022}
P.~Scholze,
``Liquid tensor experiment,''
\emph{Experimental Mathematics}, 2022.

\bibitem{LoefflerStoll2025}
D.~Loeffler and M.~Stoll,
``Formalising $\zeta$- and $L$-functions in Lean~4,''
\emph{Annals of Formalized Mathematics}, vol.~1, 2025.

\bibitem{OEIS_A003273}
OEIS Foundation,
``A003273: Congruent numbers,''
\emph{The On-Line Encyclopedia of Integer Sequences},
\url{https://oeis.org/A003273}.

\bibitem{FLTregular2025}
R.~Brasca, E.~Riccardo, et al.,
``FLT-regular: A formalization of the regular-prime case of
Fermat's Last Theorem in Lean~4,''
\emph{Annals of Formalized Mathematics} (and project repository
\url{https://github.com/leanprover-community/flt-regular}), 2025.

\end{thebibliography}

% ═════════════════════════════════════════════════════════════════════
\section*{Computational and AI-Assistance Disclosure}

During the preparation of this work the author used GitHub Copilot
(Microsoft) and Anthropic Claude in order to assist with code
generation, numerical computations, and manuscript drafting.  After
using these tools, the author reviewed and edited the content as
needed and takes full responsibility for the content of this
article.

The computational experiments, Lean development, and auxiliary
numerical checks were carried out within a reproducible, multi-tool
workflow.  All mathematical statements, Lean proofs, and numerical
results were independently verified by the author using explicit
Lean~4 proofs and reproducible Python scripts based on
\texttt{mpmath} version~1.3.0.  The full codebase, test logs, and
archived artifacts supporting this work are available at the project
repository.

Specifically, GitHub Copilot assisted with Lean~4 tactic suggestions
and boilerplate structure; Anthropic Claude assisted with
manuscript prose drafting and LaTeX formatting.  All Lean theorem
statements and proof steps were verified by the author against the
Lean~4 type checker.

\end{document}
